\subsection{Archivo de Conexión \texttt{(conexion.c)}}
Este archivo es el encargado de contener toda la lógica de la red P2P, actuando como un proceso independiente (\textit{backend}) para aislar la comunicación y evitar problemas de sincronización o bloqueos en la interfaz gráfica (GUI). Implementa una arquitectura multihilo para gestionar simultáneamente el tráfico de red mediante sockets TCP/IP y la comunicación interprocesos (IPC) mediante tuberías.

A continuación, se describe a detalle el propósito y funcionamiento de cada subrutina contenida en este módulo:

\begin{itemize}
    \item \textbf{\texttt{conexion()}:} Es la función orquestadora del \textit{backend}. Se encarga de instanciar y levantar los 4 hilos principales de ejecución mediante la directiva \texttt{pthread\_create}. Además, utiliza \texttt{pthread\_join} para mantener el proceso activo y garantizar que todos los hilos se unan y se destruyan de forma segura al momento de cerrar la aplicación.
    
    \item \textbf{\texttt{hilo\_lector\_p2p()}:} Actúa como el servidor pasivo de la arquitectura P2P. Configura un socket TCP, lo enlaza (\texttt{bind}) al puerto local y se queda a la escucha (\texttt{listen}). Cuando recibe una conexión entrante a través de \texttt{accept()}, extrae el flujo binario hacia la estructura \texttt{Mensaje} y activa la bandera volátil \texttt{regresa} para notificar que hay datos listos.
    
    \item \textbf{\texttt{hilo\_escritor\_p2p()}:} Funciona como el cliente activo de la red. Ejecuta un ciclo continuo evaluando la bandera volátil \texttt{escribe}. Cuando esta se enciende (indicando que la interfaz ha depositado un mensaje), extrae la información del buffer compartido y despacha el paquete hacia la IP y puerto del nodo remoto.
    
    \item \textbf{\texttt{hilo\_lector\_pipe()}:} Es el puente de entrada desde el \textit{frontend}. Permanece bloqueado leyendo el extremo de lectura del Pipe B (\texttt{local->B[0]}). Al recibir un paquete proveniente de la GUI, evalúa si es un texto ordinario o un evento de cierre (\texttt{id\_emisor == -1}), en cuyo caso inyecta un mensaje local para iniciar la secuencia de apagado global.
    
    \item \textbf{\texttt{hilo\_escritor\_pipe()}:} Es el puente de salida hacia la interfaz. Monitorea la bandera que levanta el servidor P2P y, al activarse, inyecta el mensaje recibido por la red en el extremo de escritura del Pipe A (\texttt{local->A[1]}) para que la vista gráfica pueda renderizarlo dinámicamente.
    
    \item \textbf{\texttt{enviar\_msj()}:} Función que encapsula la rutina de transmisión TCP. Crea un descriptor temporal, adapta la IP y el puerto a formato de red con \texttt{inet\_pton} y \texttt{htons}, establece una conexión activa de tres vías (\texttt{connect()}) con el destinatario, transmite la estructura \texttt{Mensaje} en crudo y libera el socket.
    
    \item \textbf{Funciones de Tiempo (\texttt{get\_time\_ms} y \texttt{formatear\_tiempo\_str}):} Subrutinas de apoyo encargadas de capturar la hora del sistema operativo con precisión de milisegundos y darle formato de cadena (\texttt{HH:MM:SS:ms}) para su posterior visualización en las burbujas del chat.
\end{itemize}

\begin{lstlisting}[style=estiloC, caption={Archivo del Back-End (conexion P2P)}]
#include <stdio.h>
#include <stdlib.h>
#include <string.h>
#include <unistd.h>
#include <arpa/inet.h>
#include <netinet/in.h>
#include <sys/wait.h>
#include <pthread.h>
#include <sys/time.h>
#include "header.h"

// Variable que controla el ciclo de vida de los hilos de la conexion P2P
volatile int encendido = 1;
// Variable que controla el ciclo de vida de los hilos de los pipes
volatile int activo = 1;
// Variable que controla la accion de escribir por la Red (socket)
volatile int escribe = 0;
// Variable que controla la accion de escribir para la UI (pipes)
volatile int regresa = 0;

int PUERTO = 4000;              // El puerto que usamos nosotros
char DIRECTION[] = "127.0.0.1"; // No se cambia la ip de nuestra compu

// Mensaje que se recibe por la Interfaz.
Mensaje interfaz;
// Mensaje que se recibe por el Socket.
Mensaje externo;

// Funcion que se encarga de crear los 4 hilos de la conexion P2P
void conexion(Tuberia *local)
{
    // Se crean 4 variables tipo hilo
    pthread_t servidor, escritor, pipe_lector, pipe_escritor;

    printf("[HIJO] Inicializando los 4 hilos (Sockets y Pipes)...\n");
    // Creamos el hilo servidor (P2P)
    pthread_create(&servidor, NULL, hilo_lector_p2p, NULL);
    // Creamos el hilo cliente (P2P)
    pthread_create(&escritor, NULL, hilo_escritor_p2p, NULL);
    // Creamos el hilo lector (pipe B)
    pthread_create(&pipe_lector, NULL, hilo_lector_pipe, (void *)local);
    // Creamos el hilo escritor (pipe A)
    pthread_create(&pipe_escritor, NULL, hilo_escritor_pipe, (void *)local);

    // Esperamos a los 4 hilos creados para finalizar correctamente en orden
    pthread_join(servidor, NULL);
    pthread_join(escritor, NULL);
    pthread_join(pipe_lector, NULL);
    pthread_join(pipe_escritor, NULL);

    printf("[HIJO] Todos los hilos destruidos de forma segura. Saliendo.\n");
}

// Hilo que se encarga de escuchar mensajes por el socket y nuestro PUERTO
void *hilo_lector_p2p(void *arg)
{
    // Descriptores de archivo: server_fd para la escucha y new_socket para la conexion aceptada
    int server_fd, new_socket;

    // Estructura de red que almacenara los datos de direccionamiento IP y Puerto del servidor
    struct sockaddr_in address;

    // Variable de configuracion para habilitar opciones en el nivel de socket
    int opt = 1;

    // Estructura local donde se desempaquetara el mensaje binario recibido por la red
    Mensaje msg_recibido;

    // Intentamos crear un socket de flujo (TCP) utilizando el protocolo de internet IPv4
    if ((server_fd = socket(AF_INET, SOCK_STREAM, 0)) < 0)
    {
        // Si la creacion del socket falla, terminamos el hilo de forma segura retornando NULL
        pthread_exit(NULL);
    }

    // Permite reutilizar el puerto inmediatamente si la aplicacion se reinicia, evitando el error "Address already in use"
    setsockopt(server_fd, SOL_SOCKET, SO_REUSEADDR, &opt, sizeof(opt));

    // Indicamos que utilizaremos la familia de direcciones IPv4
    address.sin_family = AF_INET;

    // Vinculamos el servidor a cualquier interfaz de red disponible en el equipo (0.0.0.0)
    address.sin_addr.s_addr = INADDR_ANY;

    // Asignamos el puerto de escucha (4000) convirtiendolo al formato de red (Big Endian) con htons
    address.sin_port = htons(PUERTO);

    // Asociamos la configuracion de la IP y el PUERTO al descriptor de archivo del socket
    if (bind(server_fd, (struct sockaddr *)&address, sizeof(address)) < 0)
    {
        // Si el puerto ya esta ocupado o falla el enlace, cerramos el descriptor y salimos del hilo
        close(server_fd);
        pthread_exit(NULL);
    }

    // Configuramos el socket en modo pasivo, permitiendo una cola de hasta 5 conexiones en espera
    if (listen(server_fd, 5) < 0)
    {
        // Si el sistema operativo no puede asignar la cola, liberamos el recurso y salimos
        close(server_fd);
        pthread_exit(NULL);
    }

    // Ciclo de vida del hilo servidor (P2P); se ejecutara activamente mientras la bandera global sea verdadera
    while (encendido)
    {
        // Obtenemos el tamaño de la estructura de direccionamiento para la llamada accept
        int tam_addr = sizeof(address);

        // Bloquea el hilo aqui hasta que un nodo remoto intente conectarse, generando un nuevo socket especifico
        new_socket = accept(server_fd, (struct sockaddr *)&address, (socklen_t *)&tam_addr);

        // Si la conexion se establecio correctamente y el descriptor es valido
        if (new_socket >= 0)
        {
            // Limpiamos la memoria del buffer local de recepcion para evitar basura de datos anteriores
            memset(&msg_recibido, 0, sizeof(Mensaje));

            // Leemos el flujo binario proveniente del socket remoto y lo mapeamos directo en la estructura Mensaje
            read(new_socket, &msg_recibido, sizeof(Mensaje));

            // Si el id_emisor es -1, significa que el nodo remoto o la interfaz local envio un evento de cierre
            if (msg_recibido.id_emisor == -1)
            {
                // Desplegamos en consola el evento global de finalizacion del programa
                printf("[GLOBAL] Evento de cierre... Cerrando hilos...\n");

                // Apagamos las banderas de control del backend para romper el ciclo de los demas hilos
                encendido = 0;
                activo = 0;

                // Cerramos de forma segura las conexiones
                close(new_socket);

                // Rompemos el ciclo infinito de escucha activa
                break;
            }
            // Si el mensaje es ordinario (contiene texto de un chat activo)
            else
            {
                // Informamos que se recibio con exito un paquete de datos en el puerto local configurado
                printf("Hilo recibio mensaje en %d...\n", PUERTO);

                // Forzamos el tipo de mensaje a RECEPCION para que la GUI sepa que debe pintarlo a la izquierda
                msg_recibido.tipo = RECEPCION;

                // Copiamos el contenido de forma segura a la variable global compartida 'externo'
                memcpy(&externo, &msg_recibido, sizeof(Mensaje));

                // Activamos la bandera volatil para indicarle al hilo escritor del pipe que hay datos listos
                regresa = 1;
            }
            // Cerramos el socket de la sesion actual para liberar el descriptor una vez procesado el mensaje
            close(new_socket);
        }
    }
    // Al salir del ciclo principal, cerramos definitivamente el socket maestro del servidor
    close(server_fd);

    // Terminamos la ejecucion del hilo de forma limpia ante el entorno de pthreads
    pthread_exit(NULL);
}

// Hilo que se encarga de escribir a un PUERTO e IP enviado por la interfaz
void *hilo_escritor_p2p(void *arg)
{
    // Ciclo de vida del hilo cliente (P2P); se ejecutara activamente mientras la bandera global sea verdadera
    while (encendido)
    {
        // Si la bandera cambia es que recibio un mensaje de la interfaz y debe enviarlo al p2p amigo
        if (escribe == 1)
        {
            // Se desactiva la bandera ya que solo debe realizarse una vez.
            escribe = 0;
            // Envia el mensaje que ya se cargo en la estructura de Mensaje interfaz;
            enviar_msj(interfaz.ip_destino, interfaz.puerto_destino, interfaz);
            // Imprime en pantalla que envio un mensaje (Depuracion)
            printf("Hilo envio mensaje a %d...\n", interfaz.puerto_destino);
        }
        // Al ser un evento volatil se debe dormir considerablemente para evitar el consumo excesivo del CPU
        usleep(1000);
    }
    // Sale del hilo para terminar la ejecucion.
    pthread_exit(NULL);
}

// Hilo que se encarga de interceptar el mensaje que envia la interfaz grafica por el pipe B
void *hilo_lector_pipe(void *arg)
{
    // El hilo tiene como entrada la estructura de Tuberia que tiene los pipes a utilizar
    Tuberia *local = (Tuberia *)arg;
    // Variable para almacenar el mensaje y despues copiarlo a la variable volatil interfaz
    Mensaje msg_pipe;

    // Ciclo de vida del hilo lector (pipe B); Se ejecuta activamente mientras la bandera global sea verdadera.
    while (activo)
    {
        // Al ser un evento bloqueante se queda esperando hasta recibir algo de la interfaz
        int bytes = read(local->B[0], &msg_pipe, sizeof(Mensaje));
        // Si el mensaje si contiene algo y sigue activo la bandera global evalua el mensaje recibido.
        if (bytes > 0 && activo)
        {
            // Si el id del mensaje es -1 significa que es un eveno de cierre y debe avisar a la conexion P2P
            if (msg_pipe.id_emisor == -1)
            {
                // Inyecta el mensaje a si mismo para que lo procese
                enviar_msj(DIRECTION, PUERTO, msg_pipe);
                // Rompemos el ciclo ya que se inyecto el mensaje de cierre y la conexion cierra los demas hilos
                break;
            }
            else // Si el id no es -1 es del usuario y debemos activar al cliente de la conexion P2P
            {
                // Muestra un mensaje en la interfaz de que se recibio un mensaje del usuario
                printf("Recibi Mensaje de la GUI...\n");
                // Copiamos el mensaje recibido a la variable global interfaz
                memcpy(&interfaz, &msg_pipe, sizeof(Mensaje));
                // Acivamos la bandera de escritura para el P2P (cliente)
                escribe = 1;
            }
        }
    }
    pthread_exit(NULL);
}

// Hilo que se encarga de escribir hacia la interfaz grafica por el pipe A
void *hilo_escritor_pipe(void *arg)
{
    // El hilo tiene como entrada la estructura de Tuberia que tiene los pipes a utilizar
    Tuberia *local = (Tuberia *)arg;
    // Ciclo de vida del hilo escritor (pipe A); Se ejecuta activamente mientras la bandera global sea verdadera.
    while (activo)
    {
        // Si la bandera cambio significa que la conexion P2P (servidor) recibio un mensaje por lo que debemos enviarlo a la interfaz por el pipe
        if (regresa == 1)
        {
            // Desactivamos la bandera ya que solo debe realizarse una vez
            regresa = 0;
            // Escribe a la tuberia unicamente (pipe A)
            write(local->A[1], &externo, sizeof(Mensaje));
        }
        // Al ser un evento volatil se debe dormir considerablemente para evitar el consumo excesivo del CPU
        usleep(1000);
    }
    pthread_exit(NULL);
}

// Funcion para enviar el Mensaje a una IP y PUERTO especifico (Cliente TCP).
void enviar_msj(char *ip_destino, int puerto_destino, Mensaje msg)
{
    // Descriptor de archivo para el socket local que iniciará la conexión.
    int sock = 0;
    
    // Estructura de red que almacenará la dirección IP y el Puerto del nodo remoto (destino).
    struct sockaddr_in serv_addr;

    // Intentamos crear un socket de flujo activo utilizando el protocolo TCP/IP bajo IPv4.
    if ((sock = socket(AF_INET, SOCK_STREAM, 0)) < 0)
        // Si el sistema operativo no puede asignar el socket, salimos de la función inmediatamente.
        return;

    // Convertimos el puerto destino de formato de host a formato de red (Big Endian) usando htons.
    serv_addr.sin_port = htons(puerto_destino);
    
    // Especificamos que la dirección de destino pertenece a la familia IPv4.
    serv_addr.sin_family = AF_INET;

    // Convertimos la IP de texto (ej. "127.0.0.1") a su representación binaria de red.
    if (inet_pton(AF_INET, ip_destino, &serv_addr.sin_addr) <= 0)
    {
        // Si la IP de destino es inválida o el formato es incorrecto, cerramos el socket y salimos.
        close(sock);
        return;
    }

    // Intentamos establecer una conexión activa de tres vías (Three-way handshake) con el nodo remoto.
    if (connect(sock, (struct sockaddr *)&serv_addr, sizeof(serv_addr)) >= 0)
    {
        // Si la conexión es exitosa, enviamos la estructura completa 'Mensaje' de manera directa por el flujo del socket.
        send(sock, &msg, sizeof(Mensaje), 0);
    }

    // Cerramos el socket una vez terminada la transferencia (o si falló la conexión) para liberar el descriptor.
    close(sock);
}

// Función para obtener el tiempo en milisegundos
long long get_time_ms()
{
    struct timeval tv;
    gettimeofday(&tv, NULL);
    return (long long)(tv.tv_sec) * 1000 + (tv.tv_usec) / 1000;
}

// Función para transformar ms a HH:MM:SS:ms
void formatear_tiempo_str(long long tiempo_ms, char *buffer_salida)
{
    time_t segundos = tiempo_ms / 1000;
    int milisegundos = tiempo_ms % 1000;
    struct tm *tm_info = localtime(&segundos);

    sprintf(buffer_salida, "%02d:%02d:%02d:%03d", tm_info->tm_hour, tm_info->tm_min, tm_info->tm_sec, milisegundos);
}
\end{lstlisting}